% Intended LaTeX compiler: pdflatex
\documentclass[a4paper,12pt]{article}
\usepackage[utf8]{inputenc}
\usepackage[T1]{fontenc}
\usepackage{graphicx}
\usepackage{longtable}
\usepackage{wrapfig}
\usepackage{rotating}
\usepackage[normalem]{ulem}
\usepackage{amsmath}
\usepackage{amssymb}
\usepackage{capt-of}
\usepackage{hyperref}
\usepackage[margin=0.2in]{geometry}
\usepackage{amsmath}
\usepackage{amsfonts}
\usepackage{indentfirst}
\usepackage{pgffor}
\usepackage{amssymb}
\usepackage{cancel}
\usepackage{amsthm}
\usepackage{polynom}
\usepackage{tikz}
\usepackage[tikz]{bclogo}
\usepackage{listings}
\usepackage[most]{tcolorbox}
\usepackage{forest}
\usepackage{adjustbox}
\usepackage{tikz-3dplot}
\usepackage{mathtools}
\usepackage{pgfplots}
\usetikzlibrary{tikzmark,calc,fit,matrix,arrows,automata,positioning}
\usepackage{centernot}
\usepackage{lmodern}
\usepackage{physics}
\usepackage[boxed,linesnumbered,vlined]{algorithm2e}
\usepackage{algpseudocode}
\usepackage{algorithmicx}
\usepackage{svg}
\tikzstyle{every picture}+=[remember picture]
\DeclareMathOperator{\spn}{span}
\DeclareMathOperator{\dom}{domain}
\DeclareMathOperator{\ran}{range}
\DeclareMathOperator{\rng}{range}
\DeclareMathOperator{\img}{Im}
\DeclareMathOperator{\adj}{adj}
\newcommand\dif[1]{\:\textrm{d}#1}
\DeclarePairedDelimiter\ceil{\lceil}{\rceil}
\DeclarePairedDelimiter\floor{\lfloor}{\rfloor}
\newcommand{\ihat}{\hat{\textbf{\i}}}
\newcommand{\jhat}{\hat{\textbf{\j}}}
\newcommand{\khat}{\hat{\textbf{k}}}
\newcommand{\rhat}{\hat{\textbf{r}}}
\newcommand{\thetahat}{\boldsymbol{\hat{\theta}}}
\author{mahmood}
\date{\today}
\title{algorithms homework 4}
\hypersetup{
 pdfauthor={mahmood},
 pdftitle={algorithms homework 4},
 pdfkeywords={},
 pdfsubject={homework on },
 pdfcreator={Emacs 30.0.50 (Org mode 9.6)}, 
 pdflang={English}}
\begin{document}

\maketitle
\definecolor{Brown}{cmyk}{0,0.81,1,0.60}
\definecolor{OliveGreen}{cmyk}{0.64,0,0.95,0.40}
\definecolor{CadetBlue}{cmyk}{0.62,0.57,0.23,0}
\definecolor{lightlightgray}{gray}{0.9}
\lstset{
  language=C,                             % Code langugage
  basicstyle=\ttfamily,                   % Code font, Examples: \footnotesize, \ttfamily
  keywordstyle=\color{OliveGreen},        % Keywords font ('*' = uppercase)
  commentstyle=\color{gray},              % Comments font
  numbers=left,                           % Line nums position
  numberstyle=\tiny,                      % Line-numbers fonts
  stepnumber=1,                           % Step between two line-numbers
  numbersep=5pt,                          % How far are line-numbers from code
  backgroundcolor=\color{lightlightgray}, % Choose background color
  frame=single,                           % A frame around the code
  tabsize=2,                              % Default tab size
  captionpos=b,                           % Caption-position = bottom
  breaklines=true,                        % Automatic line breaking?
  breakatwhitespace=false,                % Automatic breaks only at whitespace?
  showspaces=false,                       % Dont make spaces visible
  showtabs=false,                         % Dont make tabls visible
  columns=flexible,                       % Column format
  morekeywords={__global__, __device__},  % CUDA specific keywords
  showstringspaces=false,                 % dont show spaces in strings
}
% so that latex doesnt complain about sage not being a language
\lstdefinelanguage{sage}{}

% environment definition for special blocks
\theoremstyle{definition}
\newtheorem{my_example}{Example}
\counterwithin{my_example}{section}
\counterwithin{my_example}{subsection}
\newtheorem{definition}{Definition}
\counterwithin{definition}{section}
\counterwithin{definition}{subsection}
\newtheorem{characteristic}{Characteristic}
\newtheorem{entailment}{Entailment}
%\newtheore*{proof}{Proof}
\newtheorem{lemma}{Lemma}
\newtheorem{question}{Question}
\newtheorem{subquestion}{Subquestion}[question]
%\newtheorem{note}{Note}
\newtheorem{answer}{Answer}
\counterwithin{answer}{question}
\counterwithin{answer}{subquestion}
\newtheorem{step}{Step}[answer]

% the following snippet auto resizes images to fit properly
% Determine if the image is too wide for the page.
% \makeatletter
% \def\ScaleIfNeeded{%
%   \ifdim\Gin@nat@width>\linewidth
%     \linewidth
%   \else
%     \Gin@nat@width
%   \fi
% }
% \makeatother
% % Resize figures that are too wide for the page.
% \let\oldincludegraphics\includegraphics
% \renewcommand\includegraphics[2][]{%
%   \oldincludegraphics[width=\ScaleIfNeeded]{#2}
% }

\newenvironment{note}
  {\begin{bclogo}[logo=\bcinfo, couleurBarre=orange, couleur=lightgray]{ note}}
  {\end{bclogo}}
\newenvironment{attention}
  {\begin{bclogo}[logo=\bcattention, couleurBarre=red, noborder=true, 
                couleur=red!15]{Important!}}
  {\end{bclogo}}

\begin{question}
given a \href{2022-12-14T21_56_40.428Z.Desktop.org}{binary tree} \texttt{T} in which each \href{2022-12-14T21_56_40.428Z.Desktop.org}{node} has a number stored in \texttt{node.info} and a pointer to the left subtree \texttt{node.left} and the right subtree \texttt{node.right}, write an algorithm to find the average of the numbers in the tree\\[0pt]
\begin{answer}
\begin{center}
\includesvg[width=.9\linewidth]{/Users/mahmooz/.emacs.d/latex/FHjzayv}
\end{center}\\[0pt]
\end{answer}
\begin{subquestion}
prove the correctness of the algorithm\\[0pt]
\begin{answer}
\uline{proof that the algorithm terminates}:\\[0pt]
assume that \(\textsc{Average}\) runs on a tree \texttt{T}, on each recurred call, the input node is in a deeper level on the tree, so its obvious that at some point we would reach the last level in the tree and the function wouldnt recur further\\[0pt]
\uline{proof of correctness}:\\[0pt]
we prove this by \href{20220707193301-mathematical_induction.org}{perfect induction} on the height of the tree\\[0pt]
\begin{enumerate}
\item for a tree of height 1 we return \(\frac{total}{cnt} = \frac{node.info}{1} = node.info\) which is the correct average of the single node in the tree\\[0pt]
\item assume for every tree with height less than \texttt{n} the algorithm returns the correct average of the nodes in the tree\\[0pt]
\item for a tree with height \texttt{n}, \texttt{cnt} would be initialized to 1 and \texttt{total} to \texttt{node.info}, if there is a left subtree, we increment \texttt{cnt} by 1 and \texttt{total} by the sum of the nodes in the left subtree, as we know both the left and right subtrees have smaller heights than \texttt{n} we know the algorithm returns the correct sum for these inputs because of the \textbf{induction hypothesis}, and then the algorithm increments \texttt{total} by the sum of the right subtree and \texttt{cnt} again by 1, eventually we return \(\frac{total}{cnt}\) which would equal the correct average of the 1, 2, or 3 values we would have\\[0pt]
\end{enumerate}
\end{answer}
\end{subquestion}
\begin{subquestion}
analyze the \href{20221130014441-time_complexity.org}{time complexity} of the algorithm\\[0pt]
\begin{answer}
as far as i can tell, analyzing the time complexity with height as the variable using the recurrence \(T(h)=2T(h-1)+C\) wont get us a tight bound on the time, as the algorithm acts on each node in the tree only once the time complexity is probably \(\Theta(n)\) where \texttt{n} is just the number of nodes in the tree\\[0pt]
although after putting some thought into it, we could arrive at big theta using that recurrence because \(O\left(2^h\right) = O(n)\) and with \(\Omega(h)\) as the lower bound, both the upper and lower bounds are \href{2022-12-14T13_06_19.949Z.Desktop.org}{linear} so big theta would be linear, but im not going this route\\[0pt]
each recurred call is passed \texttt{n} nodes, minus one, halved, because we'd be removing the head, and ignoring half of the tree, minus one doesnt matter so we consider the \href{2022-12-14T13_06_16.539Z.Desktop.org}{recurrence formula}\\[0pt]
\[
T(n)=2T\left(\frac{n}{2}\right)+C
\]\\[0pt]
we found a closed form for this recurrence in the previous \href{2022-12-14T16_14_04.162Z.Desktop.org}{algorithms homework 3}:\\[0pt]
\begin{align*}
  T(n) &= 2T\left(\frac{n}{2}\right) + C\\
  T\left(\frac{n}{2}\right) &= 2T\left(\frac{n}{4}\right) + C\\
  T\left(\frac{n}{4}\right) &= 2T\left(\frac{n}{8}\right) + C\\
  T\left(\frac{n}{2^{\lg n}}\right) &= 2T\left(\frac{n}{2^{\lg n+1}}\right) + C\\
  T(n) &= 2\left(2T\left(\frac{n}{4}\right) + C\right) + C = 4T\left(\frac{n}{4}\right) + 3C\\
  &= 4\left(2T\left(\frac{n}{8}\right)+C\right)+3C = 8T\left(\frac{n}{8}\right)+7C\\
  &= 8\left(2T\left(\frac{n}{16}\right)+C\right)+7C = 16T\left(\frac{n}{16}\right)+15C\\
  &\vdots\\
  &= 2^{\lg n}T\left(\frac{n}{2^{\lg n}}\right) + \sum_{i=0}^{\lg n} 2^{i}C\\
  &= 2^{\lg n}C_1 + \sum_{i=0}^{\lg n} 2^{i}C\\
  &= 2^{\lg n}C_1 + \frac{1-2^{\lg(n)+1}}{-1}C\\
  &= 2^{\lg n}C_1-C+2^{\lg(n)+1}\\
  &\Rightarrow \Theta\left(2^{\lg n}\right) = \Theta(n)
\end{align*}
\end{answer}
\end{subquestion}
\end{question}

\begin{question}
write a \href{20220706211939-divide_and_conquer_algorithm.org}{divide-and-conquer algorithm} to check whether the sum of the nodes in an array is divisible by 4\\[0pt]
\begin{answer}
\begin{center}
\includesvg[width=.9\linewidth]{/Users/mahmooz/.emacs.d/latex/frGEu7Q}
\end{center}\\[0pt]
\end{answer}
\begin{subquestion}
prove the correctness of the algorithm\\[0pt]
\begin{answer}
we prove this by induction on the number of elements in the array\\[0pt]
\begin{enumerate}
\item for \(\textsc{Size}(A)=1\), the algorithm returns the only element in the array modulo 4, which would be 0 only if it is divisible which is the correct result\\[0pt]
\item the \textbf{induction hypothesis}, we use \href{20220707193301-mathematical_induction.org}{perfect induction}, we assume that the algorithm correctly computes the result for any input size that is less than \texttt{n}\\[0pt]
\item the \textbf{induction step}, let \texttt{A} be an arbitrary array of length \texttt{n}, the call to \(\textsc{Is-Divisble-By-Four}(A)\) executes the calls \(\textsc{Is-Divisble-By-Four}\left(A\left[0,\dots,\frac{\textsc{Size}(A)}{2}\right]\right)\) and \(\textsc{Is-Divisble-By-Four}\left(A\left[\frac{\textsc{Size}(A)}{2}+1,\dots,\textsc{Size}(A)\right]\right)\) and both of those calls are passed an array of size \(\frac{n}{2}\) and since it is less than \texttt{n} we know the calls return the correct result, by summing the results of those 2 calls and applying modulo 4 we get the 0 only if the sum is divisible by 4 which is the correct output\\[0pt]
\end{enumerate}
\end{answer}
\end{subquestion}
\end{question}

\begin{question}
given \texttt{T} is a binary tree, write an algorithm to find the minimum sum of the nodes along one of the routes from the root to a leaf\\[0pt]
\begin{answer}
\begin{center}
\includesvg[width=.9\linewidth]{/Users/mahmooz/.emacs.d/latex/MEkVFzd}
\end{center}\\[0pt]
\end{answer}
\begin{subquestion}
prove the correctness of the algorithm\\[0pt]
\begin{answer}
we prove this by induction on the height of the tree\\[0pt]
\begin{enumerate}
\item for height of 1 which is the \href{2022-12-14T13_06_16.539Z.Desktop.org}{base case}, the algorithm returns the value of the only node in the tree, which would indeed equal the sum of the only route which is the minimum too\\[0pt]
\item the \textbf{induction hypothesis}, we use \href{20220707193301-mathematical_induction.org}{perfect induction}, we assume that the algorithm correctly computes the result for any input size that is less than \texttt{n}\\[0pt]
\item the \textbf{induction step}, let \texttt{T} be an arbitrary tree of height \texttt{n}, the call to \(\textsc{Minimum-Route-Sum}(T)\) recurses with trees whose heights are less than \texttt{n} by 1, and so \texttt{left} and \texttt{right} would each contain the minimum sum of a route starting from the root until the leaves but we would return the smaller of both which would be the minimum route in the tree\\[0pt]
\end{enumerate}
\end{answer}
\end{subquestion}
\begin{subquestion}
analyze the \href{20221130014441-time_complexity.org}{time complexity} of the algorithm\\[0pt]
\begin{answer}
each recurred call is passed \texttt{n} nodes, minus one, halved, because we'd be removing the head, and ignoring half of the tree, minus one doesnt matter so we consider the \href{2022-12-14T13_06_16.539Z.Desktop.org}{recurrence formula}\\[0pt]
\[
T(n)=2T\left(\frac{n}{2}\right)+C
\]\\[0pt]
we solve it:\\[0pt]
\begin{align*}
  T(n) &= 2T\left(\frac{n}{2}\right) + C\\
  T\left(\frac{n}{2}\right) &= 2T\left(\frac{n}{4}\right) + C\\
  T\left(\frac{n}{4}\right) &= 2T\left(\frac{n}{8}\right) + C\\
  T\left(\frac{n}{2^{\lg n}}\right) &= 2T\left(\frac{n}{2^{\lg n+1}}\right) + C\\
  T(n) &= 2\left(2T\left(\frac{n}{4}\right) + C\right) + C = 4T\left(\frac{n}{4}\right) + 3C\\
  &= 4\left(2T\left(\frac{n}{8}\right)+C\right)+3C = 8T\left(\frac{n}{8}\right)+7C\\
  &= 8\left(2T\left(\frac{n}{16}\right)+C\right)+7C = 16T\left(\frac{n}{16}\right)+15C\\
  &\vdots\\
  &= 2^{\lg n}T\left(\frac{n}{2^{\lg n}}\right) + \sum_{i=0}^{\lg n} 2^{i}C\\
  &= 2^{\lg n}C_1 + \sum_{i=0}^{\lg n} 2^{i}C\\
  &= 2^{\lg n}C_1 + \frac{1-2^{\lg(n)+1}}{-1}C\\
  &= 2^{\lg n}C_1-C+2^{\lg(n)+1}\\
  &\Rightarrow \Theta\left(2^{\lg n}\right) = \Theta(n)
\end{align*}
\end{answer}
\end{subquestion}
\end{question}

\begin{question}
an array is \textbf{sorted-three} if all the elements that are dividable by elements in the array are placed in ascending order, with no importance to the order of elements that arent dividable by three\\[0pt]
prove that it isnt possible to transform an array into a sorted-three one using only the actions of comparison and placement in \(\Theta(n\log n)\) operations\\[0pt]
\begin{answer}
the algorithm would have to touch each element atleast ones and output a full array of size \texttt{n} so the time is bound by \(\Omega(n)\)\\[0pt]
There are \(n!\) possible permutations of the array. sorting the array is equivalent to discovering which of these permutations sorts it. moreover, the permutation must be uniquely identified by the results of comparisons, each of which gives us one bit of information.\\[0pt]
since \(\lg(n!)\) bits are required to identify the correct permutation, there must be at least \(\lg(n!)\) comparisons required for some input and using \textbf{Stirling’s Formula} we know that \(\Theta(\lg(n!))=\Theta(n\lg n)\)\\[0pt]
\end{answer}
\end{question}
\end{document}